
% =====================================================================
\section{T7: plurality forces a single quotient gate}
\label{sec:quotient}
% =====================================================================

The gating of \Cref{sec:gating} selects a walk within a graph. We now lift it:
when many parts act as one, the same theory applies to their quotient graph
(\Cref{def:quotient}), and a coherent collective admits exactly one gating
function on that quotient. The lift is a homomorphism---the four-place structure
of part, whole, gate, and invariant recurs one scale up---and its single-gate
conclusion is the structural condition for a family of parts to be one
collective rather than several.

\subsection{The collective as a quotient}

\begin{definition}[Collective; coherence]
\label{def:collective}
Let $\mathcal{P}=\{U_1,\dots,U_k\}$ be a partition of $\V$ into parts (connected
induced subgraphs). The \emph{collective} on $\mathcal{P}$ is the quotient graph
$\Quot=\Gph/\mathcal{P}$ (\Cref{def:quotient}), whose vertices are the parts and
whose edges are their mutual cuts. A gating function on $\Quot$
(\Cref{def:gate}) is a \emph{collective gate}; it biases which part acts next,
toward which neighbouring part. The collective is \emph{coherent} if it carries a
single collective gate that selects a unique next part at each branching part of
$\Quot$.
\end{definition}

\begin{remark}[Why this is the right lift]
\label{rem:lift}
By \Cref{lem:quotient-floor}, $\Quot$ is itself a contact graph: finite,
connected, with floor $\ge\floorc$. Hence every theorem proved for contact graphs
applies to $\Quot$ verbatim---T0 (its cuts are positive), T1 (its vertices, the
parts, are individuated by negation among parts), T5 (selecting a walk over parts
needs a gate). The collective is not a new kind of object; it is a contact graph
whose vertices happen to be the parts of another.
\end{remark}

\subsection{A coherent collective has exactly one gate}

\begin{theorem}[Single quotient gate]
\label{thm:single-gate}
Let $\Quot=\Gph/\mathcal{P}$ be a collective with at least one part of quotient
degree $\ge2$. Then:
\begin{enumerate}[label=\textup{(\roman*)},leftmargin=2.2em]
\item \textbf{Existence.} Selecting a walk over the parts requires a collective
gate \textup{(}T5 applied to $\Quot$, \Cref{thm:gate-necessary}\textup{)}.
\item \textbf{Uniqueness for coherence.} If the collective is coherent in the
sense that, at each branching part, the next part is determined, then the
selecting collective gate is unique on its support: any two selecting gates that
both make the collective coherent agree wherever either selects.
\end{enumerate}
A collective that admits two distinct selecting gates over the same parts is not
one coherent collective but two, partitioning into the sub-collectives on which
each gate is the selector.
\end{theorem}

\begin{proof}
(i) is \Cref{thm:gate-necessary} for the contact graph $\Quot$
(\Cref{rem:lift}). (ii) Suppose $\gate,\gate'$ are both selecting gates rendering
the collective coherent. At each branching part $U$, coherence means a unique
next part is determined; a selecting gate, by \Cref{def:gate}, assigns positive
forward weight to exactly the edge to that determined next part and zero to the
others. Hence at $U$ both $\gate$ and $\gate'$ have support exactly the single
edge to the determined successor, so $\gate(U,\cdot)$ and $\gate'(U,\cdot)$ have
the same support and both select it; they agree as selecting gates at $U$. As $U$
was an arbitrary branching part, $\gate$ and $\gate'$ agree wherever either
selects. If instead $\gate$ and $\gate'$ select \emph{different} successors at
some part, then at that part the collective has two competing determinations of
``what acts next''; the relation ``shares a selected successor'' partitions the
parts into blocks within which a single gate selects consistently, i.e. into
distinct sub-collectives---so the original was not one coherent collective but a
union of them.
\end{proof}

\begin{corollary}[Single-valuedness of a collective]
\label{cor:single-valued}
A coherent collective is single-gated: at each step exactly one part is selected
to act next. Two simultaneous selecting gates are the signature of fission into
two collectives. \textup{(}This is T5's single-step determinacy lifted to the
quotient, and is the collective analogue of a single trajectory having a single
next step.\textup{)}
\end{corollary}

\begin{proof}
Immediate from \Cref{thm:single-gate}.
\end{proof}

\subsection{The scaling homomorphism}

The constructions of the paper recur at the collective scale. We record the
correspondence as a homomorphism of four-place structures.

\begin{proposition}[Scale homomorphism]
\label{prop:homomorphism}
The assignment
\[
\bigl(\text{vertex},\ \Gph,\ \text{gate }\gate,\ \Res\bigr)
\ \longmapsto\
\bigl(\text{part},\ \Quot,\ \text{collective gate},\ \Res(\Quot)\bigr)
\]
carries each object of the single-graph theory to its counterpart in the quotient
theory, preserving the roles: the moving unit \textup{(}a step over vertices
$\mapsto$ a step over parts\textup{)}, the standing structure \textup{(}the graph
$\mapsto$ the quotient\textup{)}, the selector \textup{(}gate $\mapsto$ collective
gate\textup{)}, and the conserved invariant \textup{(}residual $\mapsto$ quotient
residual\textup{)}. Every theorem T0--T6 holds of $\Quot$ as a contact graph.
\end{proposition}

\begin{proof}
By \Cref{lem:quotient-floor} $\Quot$ is a contact graph, so T0
(\Cref{thm:floor}), T1 (\Cref{thm:negation}), T2--T3
(\Cref{thm:residual-region,thm:residual-invariant}), T4 (\Cref{thm:private}
applied to sub-collectives), T5 (\Cref{thm:gate-necessary}), and T6
(\Cref{thm:nonreturn}) all hold of $\Quot$ verbatim. The listed correspondence is
the identification of each object's role in $\Gph$ with the same role in $\Quot$,
which the definitions make a role-preserving map; it is a homomorphism of the
four-place structure (unit, whole, selector, invariant).
\end{proof}

\begin{remark}[Interpretation: collective purpose, on which nothing depends]
\label{rem:purpose}
Read the parts as agents, each running gated trajectories, and the collective
gate as the field that selects which agent acts next and toward what. Then
\Cref{cor:single-valued} says a coherent collective has a single such selecting
field, and two such fields mean two collectives---the structural form of a group
having one purpose, with competing purposes amounting to schism. The scale
homomorphism \Cref{prop:homomorphism} then reads as: \emph{part is to whole, and
gate is to invariant, exactly as agent is to collective, and purpose is to shared
invariant}. The collective's conserved residual $\Res(\Quot)$ is the invariant
from which the collective's surface activity derives, the analogue at this scale
of the private invariant of \Cref{sec:invariant}. We use only the graph
statements; the reading is orientation and is not relied upon.
\end{remark}
